\usepackage{xcolor}
\definecolor{orange1}{HTML}{F1753A}
\definecolor{orange2}{HTML}{F7A13E}
\definecolor{orange3}{HTML}{FC6300}
\definecolor{root}{HTML}{B2B2B2}
\usepackage{tikz}
\usepackage{pgfmath}
\usetikzlibrary{decorations.text, arrows.meta,calc,shadows.blur,shadings}
%\usepackage{tikz}
\usetikzlibrary{mindmap}

\tikzstyle{noeuds}=[every node/.style=concept, concept color=blue!40]
\tikzstyle{level 1 concept}+=[level distance=5cm , sibling angle = 120]
\tikzstyle{level 2 concept}+=[level distance=4cm , sibling angle = 60]

\renewcommand*\familydefault{\sfdefault} % Set font to serif family

% arctext from Andrew code with modifications:
%Variables: 1: ID, 2:Style 3:box height 4: Radious 5:start-angl 6:end-angl 7:text {format along path} 
\def\arctext[#1][#2][#3](#4)(#5)(#6)#7{

\draw[
    color=white,
    thick,
    line width=1.3pt,
    fill=#2
]
(#5:#4cm+#3) coordinate (above #1) arc (#5:#6:#4cm+#3)
-- (#6:#4) coordinate (right #1) -- (#6:#4cm-#3) coordinate (below right #1) 
arc (#6:#5:#4cm-#3) coordinate (below #1)
-- (#5:#4) coordinate (left #1) -- cycle;
\def\a#1{#4cm+#3}
\def\b#1{#4cm-#3}
\path[
    decoration={
        raise = -0.5ex, % Controls relavite text height position.
        text  along path,
        text = {#7},
        text align = center,        
    },
    decorate
    ]
    (#5:#4) arc (#5:#6:#4);
}

%arcarrow, this is mine, for beerware purpose...
%Function: Draw an arrow from arctex coordinate specific nodes to another 
%Arrow start at the start of arctext box and could be shifted to change the position
%to avoid go over another box.
%Var: 1:Start coordinate 2:End coordinate 3:angle to shift from acrtext box  
\def\arcarrow[#1](#2)(#3)[#4]{
    \draw[thick,-,>=latex,color=#1,line width=1pt,shorten >=-2pt, shorten <=-2pt] 
        let \p1 = (#2), \p2 = (#3), % To access cartesian coordinates x, and y.
            \n1 = {veclen(\x1,\y1)}, % Distance from the origin
            \n2 = {veclen(\x2,\y2)}, % Distance from the origin
            \n3 = {atan2(\y1,\x1)} % Angle where acrtext starts.
        in (\n3-#4: \n1) -- (\n3-#4: \n2); % Draw the arrow.
}

\def\arete{3} \def\epaisseur{5} \def\rayon{2}
\newcommand{\ruban}{(0,0)
++(0:0.57735*\arete-0.57735*\epaisseur+2*\rayon)
++(-30:\epaisseur-1.73205*\rayon)
arc (60:0:\rayon) -- ++(90:\epaisseur)
arc (0:60:\rayon) -- ++(150:\arete)
arc (60:120:\rayon) -- ++(210:\epaisseur)
arc (120:60:\rayon) -- cycle}








